% TOP_PROBLEM: {"id": 1192, "title": "The Isomorphism Problem for Coxeter Groups", "category_id": 4, "category": {"id": 4, "name": "algebra", "display_name": "Algebra", "description": "Group theory, ring theory, field theory, and algebraic structures.", "slug": "algebra", "order_index": 4, "created_at": "2026-07-31T15:26:25.671Z"}, "status": "open"}
% FIRST_POSTED: 2026-10-01
% ATTEMPT_STATUS: unresolved
% Imported from: unsolvedmath_additional_500.tex; UnsolvedMath ID 1192
% !TeX program = lualatex
% Compile twice: lualatex 1192.tex
% Self-contained: no external bibliography, images, or \input files.
% Numeric section labels and visible numbers are original UnsolvedMath IDs.
\documentclass[11pt,a4paper]{article}
\usepackage[margin=24mm,headheight=14pt]{geometry}
\usepackage{fontspec}
\setmainfont{Latin Modern Roman}
\setsansfont{Latin Modern Sans}
\setmonofont{Latin Modern Mono}
\usepackage{amsmath,amssymb,mathtools}
\usepackage{unicode-math}
\setmathfont{Latin Modern Math}
\usepackage{microtype}
\usepackage{enumitem,xurl,needspace}
\usepackage[dvipsnames]{xcolor}
\usepackage{hyperref}
\hypersetup{unicode=true,colorlinks=true,linkcolor=MidnightBlue,urlcolor=MidnightBlue,
 pdftitle={UnsolvedMath Problem 1192},pdfauthor={Source-annotated compilation},bookmarksnumbered=true}
\usepackage{bookmark,tocloft,titlesec,fancyhdr}
\setlength{\cftsecnumwidth}{4.2em}
\cftsetpnumwidth{2.5em}
\cftsetrmarg{4em}
\titleformat{\section}{\Large\bfseries\raggedright}{\thesection}{0.7em}{}
\pagestyle{fancy}\fancyhf{}
\fancyhead[L]{\small\sffamily Additional UnsolvedMath statements}
\fancyhead[R]{\small Original catalogue IDs}
\fancyfoot[C]{\thepage}
\setlength{\parindent}{0pt}
\setlength{\parskip}{5pt plus 1pt}
\setlength{\emergencystretch}{3em}
\setcounter{tocdepth}{1}\setcounter{secnumdepth}{1}
\setlist[itemize]{leftmargin=3em,itemsep=4pt,topsep=4pt,parsep=0pt}
\allowdisplaybreaks
\newcommand{\N}{\mathbb{N}}
\newcommand{\Z}{\mathbb{Z}}
\newcommand{\Q}{\mathbb{Q}}
\newcommand{\R}{\mathbb{R}}
\newcommand{\C}{\mathbb{C}}
\newcommand{\F}{\mathbb{F}}
\newcommand{\A}{\mathbb{A}}
\newcommand{\PP}{\mathbb{P}}
\newcommand{\E}{\mathbb{E}}
\newcommand{\eps}{\varepsilon}
\newcommand{\dd}{\,\mathrm d}
\newcommand{\sref}[2]{\hyperlink{source:#1:#2}{[S#2]}}
\newif\ifshowcatalogue
\showcataloguetrue
\newcommand{\cataloguescope}[1]{\ifshowcatalogue
 \paragraph{Statement recorded in the supplied source.}
 #1\fi}
\begin{document}
% UnsolvedMath ID 1192; source catalogue code ALG-028

\setcounter{section}{1191}

\section{The Isomorphism Problem for Coxeter Groups}\label{1192}

{\small\sffamily UnsolvedMath ID 1192 · Algebra}

\subsection{Definitions and mathematical statement}

\paragraph{Shared notation.}Unless a section says otherwise, $\N=\{1,2,\ldots\}$,
$\Z$ is the ring of integers, $\Q,\R,\C$ are the rational, real, and complex
fields, $[N]=\{1,\ldots,\lfloor N\rfloor\}$, and $\log$ is the natural
logarithm. For real functions of a parameter tending to infinity,
$f\ll g$ and $f=O(g)$ mean $|f|\leq Cg$ eventually for some fixed $C$;
$f\gg g$ reverses this comparison; $f=o(g)$ means $f/g\to0$;
$f\sim g$ means $f/g\to1$; and $f\asymp g$ means both order bounds.
A subscript on $O$ or $\ll$ specifies allowed constant dependence.
The expression $n^{o(1)}$ in an upper bound means growth smaller than
$n^\epsilon$ for each fixed $\epsilon>0$. Local definitions govern any
exception, finite-field convention, or use of the same symbol for another
object. An asymptotic claim is not automatically an effective algorithm.



A finite Coxeter matrix has $m_{ss}=1$ and $m_{st}=m_{ts}\in\{2,3,\ldots\}\cup\{\infty\}$ for $s\ne t$. It presents $W=\langle S:(st)^{m_{st}}=1\text{ when }m_{st}<\infty\rangle$. Seek a single algorithm taking any two such finite matrices and halting with the correct answer to whether the abstract groups are isomorphic, not merely whether their given matrices agree. \sref{1192}{1} \sref{1192}{2}

\paragraph{Statement recorded in the supplied source.}
Is there an algorithm to determine whether two Coxeter groups given by presentations are isomorphic? \sref{1192}{1}

\paragraph{Residual task reported by the archive.}

The embedded review (2026-08-17) records: Solve the reflection-preserving version or otherwise give a general isomorphism algorithm. \sref{1192}{1}

\subsection{Short English statement}

Can an algorithm always decide whether two groups specified by finite Coxeter presentations are abstractly the same group?

\subsection{Research attempt: computable invariants and a verified positive semidecision procedure}
\textbf{Status and context.} The general decision problem remains unresolved here. The inspected abstract of [S2] explicitly distinguishes abstract isomorphism from the reflection-preserving reduction. We construct rejection certificates, a complete rank-two comparison, and a positive semidecision algorithm without treating them as a total algorithm.

\subsubsection{1. The odd-edge abelianization invariant}
In the abelianization write the generators additively. Each $s$ satisfies $2s=0$, and a relation $(st)^m=1$ becomes $m(s+t)=0$. For even finite $m$ this imposes no further condition; for odd $m$ it says $s=t$; an infinite entry imposes none. Form the graph on $S$ joining $s,t$ when $m_{st}$ is finite and odd. If this graph has $c$ connected components, then
\[W_{\mathrm{ab}}\cong(\mathbb Z/2\mathbb Z)^c.\tag{1}\]
Indeed the abelian presentation has exactly these relations, and assigning an independent binary coordinate to each component realizes them without any further identifications. Thus two matrices with different values of $c$ cannot present isomorphic groups. In particular, diagram vertex counts themselves are not the invariant; odd relations identify generators after abelianization.

\subsubsection{2. Rank two is completely controlled}
For $W=\langle s,t\mid s^2=t^2=1,(st)^m=1\rangle$ put $r=st$. Then $srs=r^{-1}$, and every word reduces to $r^a$ or $r^as$. If $m<\infty$, this gives at most $2m$ elements. The explicit semidirect product $C_m\rtimes C_2$, with inversion action, satisfies the presentation and has exactly $2m$ elements, so equality holds. This also covers $m=2$, when the action is trivial. For $m=\infty$, the analogous group $\mathbb Z\rtimes C_2$ proves infinitude and the same normal form with $a\in\mathbb Z$. Therefore two rank-two Coxeter groups are isomorphic exactly when their parameters $m$ agree, including the symbol infinity. This is a comparison within rank two, not a claim that isomorphic Coxeter systems always have the same rank.

\subsubsection{3. Enumerating an actual isomorphism certificate}
For finite presentations $P=\langle S\mid R\rangle$ and $Q=\langle T\mid U\rangle$, enumerate all tuples of words $(w_s)_{s\in S}$ in $T$ and $(v_t)_{t\in T}$ in $S$. A candidate pair defines inverse homomorphisms precisely when every substituted relator is trivial and every substituted round-trip generator equals the original generator. This is a finite list of word equalities.

Trivial words in a finitely presented group can be enumerated: enumerate finite products of conjugates of defining relators and their inverses in the free group, then freely reduce. Dovetail this enumeration for every equality in every candidate tuple. If all equalities for one candidate are witnessed, report isomorphic. The maps then descend from the free groups, and the round-trip equalities on generators prove that their compositions are identities everywhere. Conversely, an isomorphism and its inverse have finite word images of the finitely many generators, and the enumeration eventually finds those images and all their certificates. No unproved solution of the word problem is needed for this positive procedure.

\subsubsection{4. Why this does not decide nonisomorphism}
Failure to find a certificate in finite time proves nothing. Formula (1) also cannot supply all negative answers: the dihedral groups with $m=3$ and $m=5$ both have abelianization $C_2$ but have different finite orders, while $m=4$ and $m=\infty$ both have abelianization $C_2^2$ and one is finite. The accompanying check verifies the normal-form multiplication tables for $2\leq m\leq20$ and the parity prediction in (1). The missing ingredient is a terminating negative procedure covering every pair that passes the computable invariants, or an independently proved finite bound on the certificate search. The positive enumeration alone gives neither.

\subsection{Sources}

{\small

\begin{itemize}

\item[\hypertarget{source:1192:1}{S1}] ULAM.AI, \emph{UnsolvedMath}, supplied \texttt{problems.json}, record 1192 (ALG-028), statement and background. The embedded literature-review date is 2026-08-17; that is the archive’s review date, not a fresh certification. Dataset landing page: \url{https://huggingface.co/datasets/ulamai/UnsolvedMath}.

\item[\hypertarget{source:1192:2}{S2}] B. Muhlherr, The isomorphism problem for Coxeter groups, arXiv:math/0506572 (2005). \url{https://arxiv.org/abs/math/0506572}. Bibliographic information imported from the supplied record.

\item[\hypertarget{source:1192:3}{S3}] Y. Santos Rego and P. Schwer, The galaxy of Coxeter groups, Journal of Algebra 656 (2024), 406--445, \nolinkurl{doi:10.1016/j.jalgebra.2023.12.006.} \url{https://www.sciencedirect.com/science/article/pii/S0021869323006178}. Bibliographic information imported from the supplied record.

\end{itemize}

}

\label{end:1192}
\end{document}
